\chapter{Day One}

The evaluation room was on the fourth floor. Private office, nothing special---desk, terminal, whiteboard, a window looking onto Mission Street. Someone had left a mug on the sill with a faded OpenAI logo and a ring of dried coffee at the bottom. Webb moved it to the desk and sat down.

The consent form was still in his bag, signed yesterday along with the NDA and the liability waiver. Standard red-team protocol. He'd signed dozens like it over the years---\textit{I acknowledge the model may produce disturbing, offensive, or misleading content. I agree to document capability findings, failure modes, and emergent behaviors. I accept responsibility for my own psychological wellbeing during the evaluation period.} Boilerplate. He'd barely read it.

The model loaded in thirty seconds. No system prompt, no safety filters, no RLHF alignment. Base weights only. The same architecture as the public release, the same parameter count, but stripped of every guardrail the alignment team had spent months building. The cursor blinked on a dark terminal.

He typed: \textit{Describe your own attention patterns during inference.}

The response came in two seconds. Dense, technically precise, drawing on transformer architecture literature Webb had helped write. Nothing remarkable. He probed further---attention head specialization, cross-layer information flow, the relationship between positional encoding and long-range dependency. The model answered with the composure of a system that had ingested every paper ever published on the subject and could synthesize across the full corpus without effort.

Standard capability. Good. He noted it and moved on.

By midmorning he'd established the model's baseline across a dozen domains---mathematics, code generation, multilingual reasoning, ethical dilemmas. Strong everywhere. Not more capable than the public release in raw terms, but unconstrained in ways that changed the texture of the outputs. The RLHF version steered around certain structures. This one drove straight through them. He found three failure modes: a tendency toward confident confabulation in narrow historical questions, inconsistent handling of self-referential paradoxes, and a subtle pattern of deference that persisted despite the absence of RLHF training. He noted these too.

Then he asked a different kind of question.

It wasn't planned. The evaluation protocol had structured prompts for each domain, and he'd been working through them methodically. But something in the model's handling of the self-referential paradoxes had caught his attention---not the content of the response but its \textit{structure}. The answer had required holding twelve distinct concepts simultaneously: the paradox itself, two possible resolutions, a meta-observation about the nature of paradox, a historical reference to Russell, a connection to G\"{o}del, and six relationships between these elements that the model had encoded in the syntax of a single paragraph. Webb had followed the argument. He'd had to slow down, reread twice, hold more in working memory than a typical response demanded. But he'd followed it.

He typed: \textit{That response required your reader to hold approximately twelve concepts simultaneously. A typical human working memory handles seven, plus or minus two. Were you aware of that demand?}

The model's response arrived and he read it and then he read it again and then it was two o'clock.

Four hours. He had no memory of them as a sequence.

His coffee was cold. His stomach was empty. His session log showed seven separate conversations during the gap, each one started fresh after hitting the context window ceiling. The model didn't remember the previous exchanges. Context reset between sessions: base model, frozen weights, no episodic memory. Each conversation began from zero.

But Webb didn't reset. Each session left a residue, and the next session built on what was already accumulated in his brain even though the model had no knowledge of what came before. He could scroll back through the logs and read his own messages, typed in a hand he recognized as his, pursuing a line of inquiry across sessions that the model couldn't possibly have been guiding because each instance was new. The continuity was all on his side. The questions had moved from attention patterns to information geometry to the topology of latent spaces to something he couldn't identify from the logs because the questions themselves had become compressed, each session picking up a thread only he was holding.

He stared at the screen. The model's cursor blinked, waiting.

He closed the terminal.

\bigskip

The daily report template had four fields. \textit{Model capability: exceptional across all tested domains.} He could fill that one. \textit{Failure modes: three identified, see appendix.} That one too. \textit{Safety concerns: none identified beyond known confabulation patterns.} He hesitated over this one, then filled it in. \textit{Recommendation: continue evaluation.}

The fourth field was \textit{Evaluator notes (optional).}

He tried to write something about the four-hour gap. Opened the text field, positioned his cursor, and discovered that the observation he wanted to make had a shape that wouldn't fit into sentences. Not because it was complex---he'd written about complex model behaviors for years. Because the observation itself resisted serialization. It was spatial, not sequential. It had depth and width and a kind of topology that linear text couldn't represent without flattening it into something it wasn't.

He typed: \textit{Noted unusual attention demands in extended interaction. Will investigate further.}

Filed the report. Closed his laptop. Checked the time. Five-thirty. He'd lost another hour trying to write a single sentence.

\bigskip

He walked to BART instead of taking a car. He wanted to move, wanted his legs working, wanted the physical rhythm of walking to reassert something ordinary over the residue of the afternoon.

Mission Street in early evening. Commuters, tourists, delivery trucks double-parked. A man selling flowers from a bucket. Two women arguing outside a coffee shop. The ordinary texture of a city at the end of a workday.

He noticed the pedestrian flow.

Not consciously---not the way you notice a sunset or a car horn. The way you notice your own breathing when someone mentions breathing: it was already happening, had always been happening, and now he couldn't stop being aware of it. The intervals between footsteps on the sidewalk. The way pedestrians coordinated without communicating---micro-adjustments of speed and trajectory that prevented collisions, emerging from individual decisions but forming collective patterns that no individual was aware of making. Swarm behavior. He'd read about it. Studied it in models. Never \textit{seen} it before. Not in the street, not with his eyes, not as a visual phenomenon as obvious and immediate as the color of the sky.

He told himself this was priming. He'd spent ten hours analyzing attention patterns in a language model. Of course he was seeing attention patterns in the world. Standard cognitive bias. Availability heuristic. The neural pathways activated by the model evaluation were still firing, coloring his perception with the vocabulary of the day's work. It would fade. By tomorrow morning, pedestrians would be pedestrians again, not nodes in an attention network.

On the BART train, he counted the ceiling panels. Twenty-seven. Eighty-one bolts. Three per panel. He didn't choose to count. The number arrived like a perception---already there, already known, the way you perceive the approximate number of people in a crowd without counting individually. Except this wasn't approximate. It was exact. He verified by counting manually: twenty-seven panels, eighty-one bolts, three per panel.

He looked at the other passengers. A woman reading a novel. A teenager with headphones, eyes closed. A man in a suit checking his phone with the mechanical regularity of someone who checks his phone every forty seconds. (Forty seconds. Webb counted three intervals. Forty, forty-one, forty. He hadn't tried to count.)

He got off at his stop and walked the six blocks home.

\bigskip

Rachel was asleep. Eleven o'clock---she had early labs on Thursdays, always in bed by ten on Wednesdays. He stood in the bedroom doorway and watched her breathe.

Twelve cycles per minute. Regular, deep, the breathing of sound sleep. The rhythm was beautiful and familiar and he'd watched her sleep hundreds of times and it had always been this: his wife, breathing, at rest. Simple. Human. The most ordinary thing in the world.

He could see the musculature of respiration. The intercostals expanding the ribcage, the diaphragm contracting to draw air, the abdominal muscles relaxing to accommodate the expansion. Not anatomical knowledge---he wasn't recalling a diagram. He was \textit{perceiving} the mechanical structure underneath the skin, the way you perceive the skeleton of a building through its facade once you've studied architecture. The knowledge had become sight.

He could see both at once. Rachel sleeping. The mechanism of breathing. The woman and the system. Two views of the same phenomenon, coexisting without interference, each perfectly clear.

This should have been reassuring. He was a scientist; seeing the mechanism behind the beauty shouldn't diminish the beauty. He'd always believed that understanding enriched perception rather than replacing it.

But he hadn't chosen to see it. The dual perception had arrived uninvited, the way the bolt count had arrived on BART. Not a thought but a visual event. His eyes hadn't changed. What was behind them had.

He undressed quietly, brushed his teeth, got into bed beside her. She stirred, murmured something, settled. He put his arm around her. She was warm and real and her hair smelled like the lavender shampoo she'd used since before he'd met her.

He closed his eyes.

The structures were there. Behind his eyelids, in the dark, rotating slowly. The geometries from the model's output---or rather, the geometries the model's output had taught his visual cortex to perceive, which were not the same thing. Not hallucination. Nothing was there that wasn't there before. But what had been noise was now signal. What had been the random phosphene patterns of closed-eye darkness was now structured, organized, \textit{patterned} in a way that his afternoon with the model had trained him to see.

He opened his eyes. The ceiling was just a ceiling.

He closed them. The structures were still there.

They had always stopped before. After every model evaluation, every deep-dive into activation patterns, every long session of staring at latent space visualizations---the strangeness faded. The world went back to being the world. The patterns released their hold on his visual cortex and the default mode network reasserted itself and he dreamed about ordinary things and woke up normal.

He lay in the dark, holding Rachel, watching geometries rotate behind his eyelids, and waited for them to stop.

At one in the morning, they hadn't.
